\documentclass[a4paper]{article}

% \usepackage[default]{fontsetup}

\usepackage{fancyhdr}
\usepackage{extramarks}
\usepackage{amsmath}
\usepackage{amsthm}
\usepackage{amsfonts}
\usepackage{amssymb}
\usepackage{tikz}
\usepackage[plain]{algorithm}
\usepackage{algpseudocode}
\usepackage{enumerate}
\usepackage{tikz}

\usetikzlibrary{automata,positioning}

%
% Basic Document Settings
%  

\topmargin=-0.2in
\evensidemargin=0in
\oddsidemargin=0in
\textwidth=6.5in
\textheight=9.5in
\headsep=0.25in

\linespread{1.1}

\pagestyle{fancy}
\lhead{\hmwkAuthorName}
\chead{\hmwkClass : \hmwkTitle}
\rhead{\firstxmark}
\lfoot{\lastxmark}
\cfoot{\thepage}

\renewcommand\headrulewidth{0.4pt}
\renewcommand\footrulewidth{0.4pt}

\setlength\parindent{0pt}

%
% Create Problem Sections
%

\newcommand{\enterProblemHeader}[1]{
    \nobreak\extramarks{}{Problem \arabic{#1} continued on next page\ldots}\nobreak{}
    \nobreak\extramarks{Problem \arabic{#1} (continued)}{Problem \arabic{#1} continued on next page\ldots}\nobreak{}
}

\newcommand{\exitProblemHeader}[1]{
    \nobreak\extramarks{Problem \arabic{#1} (continued)}{Problem \arabic{#1} continued on next page\ldots}\nobreak{}
    \stepcounter{#1}
    \nobreak\extramarks{Problem \arabic{#1}}{}\nobreak{}
}

\newcommand*\circled[1]{\tikz[baseline=(char.base)]{
		\node[shape=circle,draw,inner sep=2pt] (char) {#1};}}


\setcounter{secnumdepth}{0}
\newcounter{partCounter}
\newcounter{homeworkProblemCounter}
\setcounter{homeworkProblemCounter}{1}
\nobreak\extramarks{Problem \arabic{homeworkProblemCounter}}{}\nobreak{}

%
% Homework Problem Environment
%
% This environment takes an optional argument. When given, it will adjust the
% problem counter. This is useful for when the problems given for your
% assignment aren't sequential. See the last 3 problems of this template for an
% example.
%

\newenvironment{homeworkProblem}[2][-1]{%
    \ifnum#1>0\relax
        \setcounter{homeworkProblemCounter}{#1}
    \fi
    \section{Problem \arabic{homeworkProblemCounter} (#2 points)}
    \setcounter{partCounter}{1}
    \enterProblemHeader{homeworkProblemCounter}
}{%
    \exitProblemHeader{homeworkProblemCounter}
}

%
% Homework Details
%   - Title
%   - Class
%   - Due date
%   - Name
%   - Student ID

\newcommand{\hmwkTitle}{Homework\ \#05}
\newcommand{\hmwkClass}{Probability \& Statistics for EECS}
\newcommand{\hmwkDueDate}{2026-10-25}
\newcommand{\hmwkAuthorName}{}
\newcommand{\hmwkAuthorID}{}


%
% Title Page
%

\title{
    \vspace{2in}
    \textmd{\textbf{\hmwkClass:\\  \hmwkTitle}}\\
    \normalsize\vspace{0.1in}\small{Due\ on\ \hmwkDueDate\ at 23:59}\\
	\vspace{4in}
}

\author{
	Name: \textbf{\hmwkAuthorName} \\
	Student ID: \hmwkAuthorID}
\date{}

\renewcommand{\part}[1]{\textbf{\large Part \Alph{partCounter}}\stepcounter{partCounter}\\}

%
% Various Helper Commands
%

% Useful for algorithms
\newcommand{\alg}[1]{\textsc{\bfseries \footnotesize #1}}
% For derivatives
\newcommand{\deriv}[1]{\frac{\mathrm{d}}{\mathrm{d}x} (#1)}
% For partial derivatives
\newcommand{\pderiv}[2]{\frac{\partial}{\partial #1} (#2)}
% Integral dx
\newcommand{\dx}{\mathrm{d}x}
% Alias for the Solution section header
\newcommand{\solution}{\textbf{\large Solution}}
% Probability commands: Expectation, Variance, Covariance, Bias
\newcommand{\E}{\mathrm{E}}
\newcommand{\Var}{\mathrm{Var}}
\newcommand{\Cov}{\mathrm{Cov}}
\newcommand{\Bias}{\mathrm{Bias}}

\begin{document}


% \maketitle
% \thispagestyle{empty}
% \pagebreak

\date{
Due on Oct. 25, 2026, 23:59 UTC+8}
\title{SI 140A.02  Probability \& Statistics for EECS, Fall 2026 \\
Homework 5}
\maketitle
Read all the instructions below carefully before you start working on the assignment, and before you make a submission.
\begin{itemize}
    \item You are required to submit your homework in PDF format generated by \textbf{LaTeX}, otherwise you will get no points of this homework.
    \item You are required to write down all the major steps towards making your conclusions; otherwise you may obtain limited points of the problem.
    \item Write your homework in English; otherwise you will get no points of this homework.
    \item Any form of plagiarism will lead to $0$ point of this homework. 
\end{itemize}
\newpage

\begin{homeworkProblem}{20}

A/B testing is a randomized experiment used to compare two
treatments, Treatment A and Treatment B. Users arrive one by one
and are assigned to one of the treatments. Each trial is classified
as either a success or a failure.

Consider the following ``stay-with-a-winner'' procedure:
\begin{enumerate}[(i)]
    \item Assign the first user to Treatment A or Treatment B
    with equal probabilities.
    \item If the trial for the $n$th user is a success, assign
    the $(n+1)$st user to the same treatment. Otherwise, switch
    to the other treatment.
\end{enumerate}

The success probabilities for Treatments A and B are $a$ and $b$,
respectively, where $a\ne b$. Given the assigned treatment, these
success probabilities do not depend on the previous outcomes.

Let $p_n$ be the probability of success on the $n$th trial, and
let $a_n$ be the probability that Treatment A is assigned on the
$n$th trial.

\begin{enumerate}[(a)]
    \item Show that
    \[
        p_n=(a-b)a_n+b,
        \qquad
        a_{n+1}=(a+b-1)a_n+1-b.
    \]

    \item Use part (a) to show that
    \[
        p_{n+1}=(a+b-1)p_n+a+b-2ab.
    \]

    \item Assuming that the limit exists, find
    \[
        \lim_{n\to\infty}p_n.
    \]
\end{enumerate}

\par\noindent\textbf{Solution.}

\begin{enumerate}[(a)]
    \item
    Let $A_n$ be the event that Treatment A is assigned on trial $n$,
    and let $S_n$ be the event that trial $n$ is a success.
    \[
        \Pr(A_n)=a_n,
        \qquad
        \Pr(A_n^c)=1-a_n.
    \]
    By the law of total probability,
    \[
    \begin{aligned}
        p_n
        &=\Pr(S_n\mid A_n)\Pr(A_n)
          +\Pr(S_n\mid A_n^c)\Pr(A_n^c)\\
        &=aa_n+b(1-a_n).
    \end{aligned}
    \]
    \[
        \to\quad p_n=(a-b)a_n+b.
    \]

    Treatment A is assigned on trial $n+1$ exactly when trial $n$
    uses A and succeeds, or uses B and fails:
    \[
        A_{n+1}
        =(A_n\cap S_n)\cup(A_n^c\cap S_n^c).
    \]
    The two events on the right are disjoint.
    \[
    \begin{aligned}
        a_{n+1}
        &=\Pr(A_n\cap S_n)+\Pr(A_n^c\cap S_n^c)\\
        &=aa_n+(1-b)(1-a_n).
    \end{aligned}
    \]
    \[
        \to\quad a_{n+1}=(a+b-1)a_n+1-b.
    \]

    \item
    \[
        \because\quad (a-b)a_n=p_n-b,
    \]
    \[
    \begin{aligned}
        \therefore\quad p_{n+1}
        &=(a-b)a_{n+1}+b\\
        &=(a-b)\bigl[(a+b-1)a_n+1-b\bigr]+b\\
        &=(a+b-1)(p_n-b)+(a-b)(1-b)+b\\
        &=(a+b-1)p_n-b(a+b-1)+(a-b)(1-b)+b.
    \end{aligned}
    \]
    \[
        -b(a+b-1)+(a-b)(1-b)+b=a+b-2ab.
    \]
    \[
        \to\quad p_{n+1}=(a+b-1)p_n+a+b-2ab.
    \]

    \item
    Let
    \[
        L=\lim_{n\to\infty}p_n.
    \]
    \[
        \because\quad \lim_{n\to\infty}p_{n+1}=L,
        \qquad
        \therefore\quad L=(a+b-1)L+a+b-2ab.
    \]
    \[
        (2-a-b)L=a+b-2ab.
    \]
    \[
        \because\quad 0\le a,b\le1,\quad a\ne b,
        \qquad
        \therefore\quad 2-a-b>0.
    \]
    \[
        \to\quad
        \lim_{n\to\infty}p_n
        =\frac{a+b-2ab}{2-a-b}.
    \]
\end{enumerate}

\end{homeworkProblem}


\newpage

\begin{homeworkProblem}{20}

A fair coin is tossed repeatedly and independently, producing a
sequence of heads and tails. Let $N$ be the number of tosses
required to observe the first occurrence of the pattern $HH$.

Find the PMF of $N$.

\medskip
\noindent\textbf{Solution.}

Let
\[
p_n=\Pr(N=n).
\]
At least two tosses are required, and
\[
p_0=p_1=0,
\qquad
p_2=\Pr(HH)=\frac14.
\]

For $n\ge3$, partition according to the initial pattern:
\begin{itemize}
    \item $T$, with probability $1/2$:
    the process restarts after one toss, and $n-1$ more tosses
    are required.
    \item $HT$, with probability $1/4$:
    the process restarts after two tosses, and $n-2$ more tosses
    are required.
    \item $HH$, with probability $1/4$:
    $N=2$, contributing zero to $p_n$ for $n\ge3$.
\end{itemize}
The remaining tosses are independent of the initial pattern.
\[
\therefore\quad
p_n=\frac12p_{n-1}+\frac14p_{n-2},
\qquad n\ge3.
\]

Define
\[
a_n=2^n p_n,\qquad n\ge1.
\]
Multiplying the recurrence by $2^n$ gives
\[
a_n
=2^{n-1}p_{n-1}+2^{n-2}p_{n-2}
=a_{n-1}+a_{n-2}.
\]
The sequence is specified by
\[
a_1=0,\qquad a_2=1,\qquad
a_n=a_{n-1}+a_{n-2}\quad(n\ge3),
\]
with first terms
\[
a_1,a_2,a_3,a_4,a_5,a_6,\ldots
=0,1,1,2,3,5,\ldots.
\]

\[
\longrightarrow\quad
\Pr(N=n)=
\begin{cases}
\dfrac{a_n}{2^n}, & n=2,3,4,\ldots,\\[2mm]
0, & \text{otherwise}.
\end{cases}
\]

\end{homeworkProblem}

\newpage

\begin{homeworkProblem}{15}

A system contains seven devices: four in Group A and three in
Group B. Each device fails independently with probability $p$,
where $0<p<1$.

Let $X$ be the number of failed devices in Group A, and let $T$
be the total number of failed devices in the system.

\begin{enumerate}[(a)]
    \item Find the PMF of $X$ and the PMF of $T$.

    \item Given that exactly three devices have failed, find
    the conditional PMF of $X$. Identify the distribution.

    \item Explain why the conditional distribution in part (b)
    does not depend on $p$.
\end{enumerate}

\medskip
\noindent\textbf{Solution.}

\begin{enumerate}[(a)]
    \item
    The device failures are independent, each with probability $p$.
    \[
    \therefore\quad
    X\sim\operatorname{Binomial}(4,p),
    \qquad
    T\sim\operatorname{Binomial}(7,p).
    \]

    Their PMFs are
    \[
    P(X=x)=
    \begin{cases}
        \displaystyle \binom{4}{x}p^x(1-p)^{4-x},
        & x=0,1,2,3,4,\\[2mm]
        0, & \text{otherwise},
    \end{cases}
    \]
    and
    \[
    P(T=t)=
    \begin{cases}
        \displaystyle \binom{7}{t}p^t(1-p)^{7-t},
        & t=0,1,\ldots,7,\\[2mm]
        0, & \text{otherwise}.
    \end{cases}
    \]

    \item
    For $x=0,1,2,3$, the event $\{X=x,T=3\}$ requires
    $x$ failures in Group A and $3-x$ failures in Group B.
    By independence,
    \[
    \begin{aligned}
    P(X=x,T=3)
    &=
    \binom{4}{x}p^x(1-p)^{4-x}
    \binom{3}{3-x}p^{3-x}(1-p)^x\\
    &=
    \binom{4}{x}\binom{3}{3-x}p^3(1-p)^4.
    \end{aligned}
    \]
    Also,
    \[
    P(T=3)=\binom{7}{3}p^3(1-p)^4.
    \]

    \[
    \begin{aligned}
    P(X=x\mid T=3)
    &=
    \frac{P(X=x,T=3)}{P(T=3)}\\[1mm]
    &=
    \frac{\binom{4}{x}\binom{3}{3-x}p^3(1-p)^4}
         {\binom{7}{3}p^3(1-p)^4}\\[1mm]
    &=
    \frac{\binom{4}{x}\binom{3}{3-x}}{\binom{7}{3}}.
    \end{aligned}
    \]

    \[
    \longrightarrow\quad
    P(X=x\mid T=3)=
    \begin{cases}
        \displaystyle
        \frac{\binom{4}{x}\binom{3}{3-x}}{35},
        & x=0,1,2,3,\\[3mm]
        0, & \text{otherwise}.
    \end{cases}
    \]

    Explicitly,
    \[
    \begin{array}{c|cccc}
        x & 0 & 1 & 2 & 3\\ \hline
        P(X=x\mid T=3)
        & \dfrac{1}{35}
        & \dfrac{12}{35}
        & \dfrac{18}{35}
        & \dfrac{4}{35}
    \end{array}
    \]

    This is a hypergeometric distribution:
    \[
    X\mid(T=3)
    \sim\operatorname{Hypergeometric}(N=7,K=4,n=3),
    \]
    corresponding to selecting $3$ devices without replacement
    from $7$ devices, of which $4$ belong to Group A.

    \item
    Every specified configuration with exactly three failed
    devices and four working devices has probability
    \[
    p^3(1-p)^4.
    \]
    Its conditional probability given $T=3$ is
    \[
    \frac{p^3(1-p)^4}
         {\binom{7}{3}p^3(1-p)^4}
    =\frac{1}{35}.
    \]

    All $\binom{7}{3}=35$ such configurations are equally likely
    after conditioning on $T=3$. The common factor
    $p^3(1-p)^4$ cancels:
    \[
    \longrightarrow\quad
    P(X=x\mid T=3)
    =
    \frac{\binom{4}{x}\binom{3}{3-x}}{\binom{7}{3}},
    \qquad x=0,1,2,3,
    \]
    which does not depend on $p$.
\end{enumerate}

\end{homeworkProblem}

\newpage

\begin{homeworkProblem}{15}

Two fair six-sided dice are rolled independently. Let $U$ and
$V$ be their respective outcomes, and define
\[
    D=|U-V|.
\]

\begin{enumerate}[(a)]
    \item Find the PMF of $D$.

    \item Find the CDF $F_D(t)$ for every real number $t$.

    \item Define $W=\min(D,2)$. Find the PMF of $W$.
\end{enumerate}

\medskip
\noindent\textbf{Solution.}

The $36$ ordered outcomes are equally likely:
\[
P(U=u,V=v)
=P(U=u)P(V=v)
=\frac{1}{36},
\qquad u,v\in\{1,\ldots,6\}.
\]

\begin{enumerate}[(a)]
    \item The possible values of $D$ are $0,1,\ldots,5$.

    For $d=0$,
    \[
    P(D=0)
    =\sum_{k=1}^{6}P(U=k,V=k)
    =\frac{6}{36}
    =\frac16.
    \]

    For $d\in\{1,\ldots,5\}$,
    \[
    \begin{aligned}
    P(D=d)
    &=P(U-V=d)+P(V-U=d)\\
    &=\sum_{k=1}^{6-d}
      \left[
      P(U=k+d,V=k)+P(U=k,V=k+d)
      \right]\\
    &=\frac{2(6-d)}{36}
    =\frac{6-d}{18}.
    \end{aligned}
    \]

    \[
    \rightarrow\quad
    p_D(d)=
    \begin{cases}
        \dfrac16, & d=0,\\[2mm]
        \dfrac{6-d}{18}, & d\in\{1,2,3,4,5\},\\[2mm]
        0, & \text{otherwise}.
    \end{cases}
    \]

    \item By definition,
    \[
    F_D(t)=P(D\le t)
    =\sum_{\substack{d\in\{0,1,\ldots,5\}\\d\le t}}p_D(d).
    \]

    The cumulative probabilities are
    \[
    \begin{aligned}
    P(D\le0)&=\frac{6}{36}=\frac16,\\
    P(D\le1)&=\frac{6+10}{36}=\frac49,\\
    P(D\le2)&=\frac{6+10+8}{36}=\frac23,\\
    P(D\le3)&=\frac{6+10+8+6}{36}=\frac56,\\
    P(D\le4)&=\frac{6+10+8+6+4}{36}=\frac{17}{18},\\
    P(D\le5)&=1.
    \end{aligned}
    \]

    \[
    \rightarrow\quad
    F_D(t)=
    \begin{cases}
        0, & t<0,\\[1mm]
        \dfrac16, & 0\le t<1,\\[2mm]
        \dfrac49, & 1\le t<2,\\[2mm]
        \dfrac23, & 2\le t<3,\\[2mm]
        \dfrac56, & 3\le t<4,\\[2mm]
        \dfrac{17}{18}, & 4\le t<5,\\[2mm]
        1, & t\ge5.
    \end{cases}
    \]

    \item The definition of $W$ gives
    \[
    W=
    \begin{cases}
        0, & D=0,\\
        1, & D=1,\\
        2, & D\ge2.
    \end{cases}
    \]

    \[
    \begin{aligned}
    P(W=0)&=P(D=0)=\frac16,\\
    P(W=1)&=P(D=1)=\frac5{18},\\
    P(W=2)&=P(D\ge2)\\
          &=1-P(D=0)-P(D=1)\\
          &=1-\frac16-\frac5{18}
          =\frac59.
    \end{aligned}
    \]

    \[
    \rightarrow\quad
    p_W(w)=
    \begin{cases}
        \dfrac16, & w=0,\\[2mm]
        \dfrac5{18}, & w=1,\\[2mm]
        \dfrac59, & w=2,\\[2mm]
        0, & \text{otherwise}.
    \end{cases}
    \]
\end{enumerate}

\end{homeworkProblem}

\newpage

\begin{homeworkProblem}{15}

There are two coins, A and B. Coin A lands heads with probability
$1/4$, while coin B lands heads with probability $3/4$.

One coin is selected with equal probabilities and then tossed
twice. The same selected coin is used for both tosses. Given
the selected coin, the two tosses are independent.

Let $C\in\{A,B\}$ denote the selected coin. Let $X$ and $Y$ be
the indicator random variables for heads on the first and
second tosses, respectively.

\begin{enumerate}[(a)]
    \item Find the PMF of $X$ and the PMF of $Y$.

    \item Are $X$ and $Y$ independent? Are they conditionally
    independent given $C$? Justify your answers.

    \item Given that both tosses are heads, find the
    conditional probability that coin B was selected.
\end{enumerate}

\medskip
\noindent\textbf{Solution.}

\begin{enumerate}[(a)]
    \item
    By the law of total probability,
    \[
    \begin{aligned}
    P(X=1)
    &=P(C=A)P(X=1\mid C=A)\\
    &\quad+P(C=B)P(X=1\mid C=B)\\
    &=\frac12\cdot\frac14+\frac12\cdot\frac34
    =\frac12.
    \end{aligned}
    \]
    Similarly,
    \[
    P(Y=1)
    =\frac12\cdot\frac14+\frac12\cdot\frac34
    =\frac12.
    \]
    Also,
    \[
    P(X=0)=1-P(X=1)=\frac12,
    \qquad
    P(Y=0)=1-P(Y=1)=\frac12.
    \]
    \[
    \Longrightarrow\quad
    p_X(k)=p_Y(k)=
    \begin{cases}
        \dfrac12, & k\in\{0,1\},\\[2mm]
        0, & \text{otherwise}.
    \end{cases}
    \]

    \item
    Using conditional independence given the selected coin,
    \[
    \begin{aligned}
    P(X=1,Y=1)
    &=P(C=A)P(X=1,Y=1\mid C=A)\\
    &\quad+P(C=B)P(X=1,Y=1\mid C=B)\\
    &=\frac12\left(\frac14\right)^2
      +\frac12\left(\frac34\right)^2\\
    &=\frac{1}{32}+\frac{9}{32}
    =\frac{5}{16}.
    \end{aligned}
    \]
    However,
    \[
    P(X=1)P(Y=1)=\frac12\cdot\frac12=\frac14.
    \]
    \[
    P(X=1,Y=1)\ne P(X=1)P(Y=1)
    \quad\Longrightarrow\quad
    X\text{ and }Y\text{ are not independent}.
    \]

    For conditional independence, write
    \[
    p_A=\frac14,\qquad p_B=\frac34.
    \]
    The problem states that the two tosses are independent
    for each fixed selected coin. For every $c\in\{A,B\}$
    and $x,y\in\{0,1\}$,
    \[
    \begin{aligned}
    P(X=x,Y=y\mid C=c)
    &=p_c^{x+y}(1-p_c)^{2-x-y}\\
    &=\bigl[p_c^x(1-p_c)^{1-x}\bigr]
      \bigl[p_c^y(1-p_c)^{1-y}\bigr]\\
    &=P(X=x\mid C=c)\,P(Y=y\mid C=c).
    \end{aligned}
    \]
    \[
    \Longrightarrow\quad
    X\text{ and }Y\text{ are conditionally independent given }C.
    \]

    \item
    By Bayes' formula,
    \[
    \begin{aligned}
    P(C=B\mid X=1,Y=1)
    &=\frac{P(C=B)\,P(X=1,Y=1\mid C=B)}
    {P(X=1,Y=1)}\\[1mm]
    &=\frac{\frac12\left(\frac34\right)^2}
    {\frac12\left(\frac14\right)^2
    +\frac12\left(\frac34\right)^2}\\[1mm]
    &=\frac{9/32}{10/32}.
    \end{aligned}
    \]
    \[
    \Longrightarrow\quad
    P(C=B\mid X=1,Y=1)=\frac9{10}.
    \]
\end{enumerate}

\end{homeworkProblem}

\newpage

\begin{homeworkProblem}{15}

A random variable $X$ has PMF
\[
    \Pr(X=1)=\frac12,\qquad
    \Pr(X=2)=\frac14,\qquad
    \Pr(X=3)=\Pr(X=4)=\frac18.
\]
Let $U$ have the discrete uniform distribution on
$\{1,2,3,4\}$.

Throughout this problem, logarithms have base 2, and entropy
is measured in bits.

\begin{enumerate}[(a)]
    \item Find $H(X)$ and $H(U)$.

    \item Define
    \[
        Y=
        \begin{cases}
            1,&X\le2,\\
            0,&X>2,
        \end{cases}
        \qquad
        Z=
        \begin{cases}
            1,&X=1,\\
            0,&X\ne1.
        \end{cases}
    \]
    Find the PMF of each variable and calculate $H(Y)$
    and $H(Z)$.

    \item Is the size of the support sufficient to determine
    the entropy of a discrete random variable?
    Explain using distributions from this problem.
\end{enumerate}

\textbf{Solution.}

\begin{enumerate}[(a)]
    \item By the definition of entropy,
    \[
        H(X)=-\sum_x \Pr(X=x)\log_2\Pr(X=x).
    \]
    \[
    \begin{aligned}
        H(X)
        &=-\left(
            \frac12\log_2\frac12
            +\frac14\log_2\frac14
            +\frac18\log_2\frac18
            +\frac18\log_2\frac18
        \right)\\
        &=\frac12+\frac12+\frac38+\frac38
        =\frac74.
    \end{aligned}
    \]
    For $U$,
    \[
        \Pr(U=u)=\frac14,\qquad u\in\{1,2,3,4\}.
    \]
    \[
        \rightarrow\quad
        H(U)
        =-\sum_{u=1}^{4}\frac14\log_2\frac14
        =2.
    \]

    \item For $Y$,
    \[
    \begin{aligned}
        \Pr(Y=1)
        &=\Pr(X\le2)
        =\frac12+\frac14
        =\frac34,\\
        \Pr(Y=0)
        &=\Pr(X>2)
        =\frac18+\frac18
        =\frac14.
    \end{aligned}
    \]
    \[
        \rightarrow\quad
        p_Y(y)=
        \begin{cases}
            \frac34,&y=1,\\[2pt]
            \frac14,&y=0,\\[2pt]
            0,&\text{otherwise}.
        \end{cases}
    \]
    \[
    \begin{aligned}
        H(Y)
        &=-\frac34\log_2\frac34
          -\frac14\log_2\frac14\\
        &=-\frac34(\log_2 3-2)+\frac12\\
        &=2-\frac34\log_2 3
        \approx0.8113.
    \end{aligned}
    \]

    For $Z$,
    \[
    \begin{aligned}
        \Pr(Z=1)&=\Pr(X=1)=\frac12,\\
        \Pr(Z=0)&=\Pr(X\ne1)=1-\frac12=\frac12.
    \end{aligned}
    \]
    \[
        \rightarrow\quad
        p_Z(z)=
        \begin{cases}
            \frac12,&z\in\{0,1\},\\[2pt]
            0,&\text{otherwise}.
        \end{cases}
    \]
    \[
        H(Z)
        =-\frac12\log_2\frac12
         -\frac12\log_2\frac12
        =1.
    \]

    \item No. The support size alone does not determine entropy.
    \[
        \operatorname{supp}(X)
        =\operatorname{supp}(U)
        =\{1,2,3,4\},
        \qquad
        H(X)=\frac74\ne2=H(U).
    \]
    Another example is
    \[
        \operatorname{supp}(Y)
        =\operatorname{supp}(Z)
        =\{0,1\},
        \qquad
        H(Y)=2-\frac34\log_2 3\ne1=H(Z).
    \]
    \[
        \rightarrow\quad
        \text{Entropy also depends on the probabilities
        assigned to the support values.}
    \]
\end{enumerate}

\end{homeworkProblem}
\end{document}